\documentclass{article}\usepackage{amsmath}\usepackage{graphicx}\usepackage{geometry}\geometry{a4paper, margin=1in}\title{SelAction Program: Technical Description of \texttt{selinbreeding.f90}}\author{Generated by Gemini}\date{\today}\begin{document}\maketitle\begin{abstract}This document provides a detailed technical description of the \texttt{selinbreeding.f90} module from the SelAction program. This module calculates the rate of inbreeding for multi-trait selection on Best Linear Unbiased Prediction (BLUP) Estimated Breeding Values (EBVs).\end{abstract}\tableofcontents\section{Rate of Inbreeding}The rate of inbreeding ($\Delta F$) is calculated using the method of Bijma and Woolliams (2000), which is based on long-term genetic contributions. The core equation is:\begin{equation}\Delta F = \frac{1}{2} \left( \frac{1}{N_m} \sigma^2_{c,m} + \frac{1}{N_f} \sigma^2_{c,f} \right)\end{equation}where $N_m$ and $N_f$ are the number of sires and dams, and $\sigma^2_{c,m}$ and $\sigma^2_{c,f}$ are the variances of long-term genetic contributions for males and females, respectively.\subsection{Long-term Genetic Contributions}The variance of long-term genetic contributions is a function of the variance of the selective advantage and the regression of an individual's breeding value on the selective advantage of its parents. The selective advantage is the sum of breeding values weighted by their economic values. The program calculates the (co)variances of EBVs and selective advantages to solve for the rate of inbreeding.\subsection{Poisson Correction}For small populations (fewer than 20 parents), a correction factor is applied to account for the fact that the distribution of family sizes is not perfectly Poisson. This correction, based on the work of Wray et al. (1990), adjusts the selected proportions to provide a more accurate prediction of the rate of inbreeding.\end{document}